% ============================================================================
%  Problem 1 - minimum Hamming distance: processing time vs. memory size
%  IEEE conference format (IEEEtran, A4)
%
%  Overleaf : New Project > Upload Project > overleaf_project.zip
%             (this file + figures/). Compiler: pdfLaTeX.
%  Figures  : figures/wave_v1.png, wave_v2.png  RTL simulation (Icarus Verilog),
%                                               made by sim/run_icarus.bat
%             figures/bd_v2.pdf                 Vivado block design (version 2;
%                                               figures/bd_v1.pdf is version 1)
%             A missing figure is printed as a framed box, so the file always
%             compiles.
% ============================================================================
\documentclass[conference,a4paper]{IEEEtran}
\usepackage[T1]{fontenc}
\usepackage{cite}
\usepackage{amsmath}
\usepackage{graphicx}
\usepackage{array,booktabs,tabularx}
\usepackage{xcolor}
\usepackage{listings}
\usepackage[scaled=0.95,varqu,varl]{inconsolata}
\usepackage{url}

\definecolor{codebg}{HTML}{F6F8FA}
\definecolor{coderule}{HTML}{D0D7DE}
\definecolor{codekw}{HTML}{1F4E79}
\colorlet{codecomment}{black!55}
\colorlet{codenumber}{black!45}

\lstset{language=C, morekeywords={size_t},
  basicstyle=\ttfamily\fontsize{6.4}{7.3}\selectfont,
  keywordstyle=\color{codekw}\bfseries,
  commentstyle=\color{codecomment},
  backgroundcolor=\color{codebg},
  frame=single, rulecolor=\color{coderule}, framerule=0.4pt, framesep=2pt,
  numbers=left, numberstyle=\tiny\color{codenumber}, numbersep=4pt,
  xleftmargin=10pt, breaklines=true,
  postbreak=\mbox{\textcolor{codenumber}{\fontsize{7}{8}\selectfont$\hookrightarrow$}\space},
  columns=fullflexible, keepspaces=true, showstringspaces=false,
  tabsize=4, upquote=true, aboveskip=4pt, belowskip=2pt}
% IEEEtran's captions add no space below them; keep the listing frame clear of
% its caption
\makeatletter
\def\lst@makecaption#1#2{\@makecaption{#1}{#2}\vskip 4pt}
\makeatother

\newcommand{\code}[1]{\texttt{#1}}
% picture figures/<name>.pdf|.png|.jpg if it exists, a framed box otherwise:
% name, width, maximum height, text of the box
\newcommand{\figorbox}[4]{%
  \IfFileExists{#1.pdf}{\includegraphics[width=#2,height=#3,keepaspectratio]{#1.pdf}}{%
  \IfFileExists{#1.png}{\includegraphics[width=#2,height=#3,keepaspectratio]{#1.png}}{%
  \IfFileExists{#1.jpg}{\includegraphics[width=#2,height=#3,keepaspectratio]{#1.jpg}}{%
  \fbox{\parbox[c][#3][c]{\dimexpr#2-2\fboxsep-2\fboxrule\relax}{\centering\small #4}}}}}}

\begin{document}

\title{Time--Memory Trade-off in Computing the Minimum Hamming Distance of
32-Bit Words: Brute Force versus Sort-and-Search}

\author{\IEEEauthorblockN{Ahmed Saad Sweffi}
\IEEEauthorblockA{Egypt-Japan University of Science and Technology (E-JUST), Alexandria, Egypt\\
ahmed.abdelgaber@ejust.edu.eg}}

\maketitle

\begin{abstract}
The task was to write a C function that returns the minimum Hamming distance
between the 32-bit words of an array, keeping both the processing time and the
memory size small. Since these goals conflict, I wrote two versions. Version~1
compares every pair of words and needs no extra memory. Version~2 sorts a copy
of the array, searches for pairs at distance 0 and 1, and only then compares
pairs, stopping at the smallest distance still possible. I ran both in C on a PC
and on a Zynq-7020, as ARM Cortex-A9 software and as Vitis HLS cores, and they
agree on every test. On 5,000 words, version~2 is 51 times faster on the PC, 37
times on the ARM and 7.8 times as an FPGA core, while version~1 needs no extra
heap, 53 times less code and data on the ARM and 4.7 times fewer LUTs.
\end{abstract}

\begin{IEEEkeywords}
Hamming distance, time--memory trade-off, sorting, binary search, popcount,
Zynq-7000, Vitis HLS, FPGA
\end{IEEEkeywords}

% ============================================================================
\section{Introduction}
The task I worked on was the following:
\begin{quote}\itshape
``Write a C function that computes the minimum Hamming distance between the
32-bit words of an array, and optimize the processing time and the memory size
required to compute the result as much as you can.''
\end{quote}
The function has a fixed prototype:
\begin{lstlisting}[numbers=none,frame=none,backgroundcolor={},xleftmargin=0pt,basicstyle=\ttfamily\footnotesize]
int min_Hamming(unsigned int tab[], unsigned int tab_size);
\end{lstlisting}
The Hamming distance between two words is the number of bit positions where
they differ, which is the number of 1-bits in their XOR~\cite{hamming}. The
function returns the smallest such distance over all pairs of words in
\code{tab}, a value from 0 to 32, or $-1$ when \code{tab} holds fewer than two
words.

The last part of the task is the hard one, because time and memory pull in
opposite directions: the fast way to find close words is to build extra data,
such as a sorted copy of the array, while the smallest program compares the
words where they are. A single version would end up somewhere in between, so I
wrote two and measured both goals for each: version~1
(\code{min\_hamming\_v1.c}) aims at the smallest memory and version~2
(\code{min\_hamming\_v2.c}) at the shortest time.

I tested both versions on two platforms: in C on a PC, and on a Zynq-7020 SoC
(XC7Z020-1), where each version runs as bare-metal software on the ARM
Cortex-A9 and as a hardware core made with Vitis HLS~\cite{ug1399} in the FPGA
fabric.

% ============================================================================
\section{The Two Versions}

\subsection{Version 1: Brute Force (Memory Size)}
Version~1 (Listing~\ref{lst:v1}) is the direct method. For every pair of words
it computes the XOR, counts the 1-bits one at a time with a shift-and-test loop,
and keeps the smallest count. It always goes through all $N(N-1)/2$ pairs, so
its time is $O(32N^2)$ whatever the input. On the other hand, it needs only a
few registers (no \code{malloc}, no library call, no copy of the data), so its
extra memory is $O(1)$.

\begin{lstlisting}[float=t,caption={min\_hamming\_v1.c (comments shortened)},label={lst:v1}]
int min_Hamming(unsigned int tab[], unsigned int tab_size) {
    if (tab_size < 2)          // no pair, no answer
        return -1;
    int best = 32;             // words differ in at most 32 bits
    for (unsigned int i = 0; i < tab_size; i++)
        for (unsigned int j = i + 1; j < tab_size; j++) {
            unsigned int diff = tab[i] ^ tab[j]; // 1 = differ
            int distance = 0;
            for (int bit = 0; bit < 32; bit++)   // bit by bit
                if ((diff >> bit) & 1u)
                    distance++;
            if (distance < best)
                best = distance;
        }
    return best;
}
\end{lstlisting}

\subsection{Version 2: Sort and Search (Processing Time)}
Version~2 (Listing~\ref{lst:v2}) is based on a simple fact: the answer is
usually small, and a small distance can be found without trying every pair. So
it checks the possible answers from the smallest one up:
\begin{enumerate}
\item \emph{Distance 0.} It sorts a copy of the array with \code{qsort}. Equal
  words end up next to each other, and one pass over the copy finds them.
  Cost: $O(N\log N)$.
\item \emph{Distance 1.} For every word, it flips each of the 32 bits in turn and
  looks for the result in the sorted copy with a binary search~\cite{knuth3}. A
  hit means distance~1. I wrote the search by hand, because \code{bsearch} would
  call the comparison function through a pointer at every step.
  Cost: $O(32N\log N)$.
\item \emph{Distance 2 or more.} It compares all pairs like version~1, but counts
  the bits with one instruction (GCC turns \code{\_\_builtin\_popcount} into
  \code{POPCNT}), and it stops at the first pair at distance~2,
  since nothing smaller is left. The worst case is still $O(N^2)$, but a pair
  now costs about 8 instructions instead of 267.
\end{enumerate}

The price is memory. The sorted copy takes $4N$ bytes (I do not sort
\code{tab[]} itself, since the caller does not expect it to change), and glibc's
\code{qsort} adds a $4N$-byte merge buffer, so the peak on the PC is $8N$ bytes;
newlib's \code{qsort} on the ARM sorts in place, so there it is $4N$. If
\code{malloc} fails, version~2 skips steps 1 and 2 and still returns the right
answer, only more slowly. Table~\ref{tab:versions} compares the two versions.

\begin{table}[!t]
\caption{The Two Versions at a Glance ($N$ Words)}
\label{tab:versions}
\centering\footnotesize\hbadness=10000
\setlength{\tabcolsep}{3pt}
\begin{tabularx}{\columnwidth}{@{}>{\raggedright}p{1.3cm}>{\raggedright}X>{\raggedright\arraybackslash}X@{}}
\toprule
 & \textbf{Version 1} & \textbf{Version 2} \\
\midrule
Idea & XOR every pair, count the bits one by one
     & sort a copy, search the one-bit neighbours, then scan the pairs with an
       early exit \\
Time & $O(32N^2)$, always all pairs
     & $O(N\log N)$ if the answer is 0, $O(32N\log N)$ if it is 1, else $O(N^2)$
       at worst \\
Extra memory & $O(1)$, no allocation
     & $O(N)$: $8N$\,B on the PC, $4N$\,B on the ARM \\
Cost per pair & 267 instructions (PC), 37 (AVX2 build), about 201 (ARM)
     & 8 instructions (PC), about 21 (ARM, which has no popcount instruction) \\
\bottomrule
\end{tabularx}
\end{table}

\begin{lstlisting}[float=t,caption={min\_hamming\_v2.c (comments shortened)},label={lst:v2}]
#include <stdlib.h>   // malloc, free, qsort
#include <string.h>   // memcpy
// one POPCNT instruction with gcc -mpopcnt (the file also has
// a portable bit count for other compilers)
#define POPCOUNT32(x) __builtin_popcount(x)
static int compare_words(const void* pa, const void* pb) {
    unsigned int a = *(const unsigned int*)pa;
    unsigned int b = *(const unsigned int*)pb;
    return (a > b) - (a < b);    // a - b could overflow
}
int min_Hamming(unsigned int tab[], unsigned int tab_size) {
    if (tab_size < 2)
        return -1;
    int lowest = 0;              // smallest answer still possible
    unsigned int* sorted =
        malloc((size_t)tab_size * sizeof(unsigned int));
    if (sorted != NULL) {        // no memory: skip steps 1 and 2
        memcpy(sorted, tab, (size_t)tab_size * sizeof(unsigned int));
        qsort(sorted, tab_size, sizeof(unsigned int), compare_words);
        // step 1: equal neighbours -> distance 0
        for (unsigned int i = 0; i + 1 < tab_size; i++)
            if (sorted[i] == sorted[i + 1]) {
                free(sorted);
                return 0;
            }
        // step 2: binary search for the 32 one-bit neighbours
        for (unsigned int i = 0; i < tab_size; i++)
            for (int bit = 0; bit < 32; bit++) {
                unsigned int neighbour = sorted[i] ^ (1u << bit);
                unsigned int lo = 0, hi = tab_size;
                while (lo < hi) {
                    unsigned int mid = lo + (hi - lo) / 2;
                    if (sorted[mid] < neighbour) lo = mid + 1;
                    else                         hi = mid;
                }
                if (lo < tab_size && sorted[lo] == neighbour) {
                    free(sorted);
                    return 1;
                }
            }
        free(sorted);
        lowest = 2;              // no pair at distance 0 or 1
    }
    // step 3: all pairs, one-instruction bit count, early exit
    int best = 32;
    for (unsigned int i = 0; i < tab_size; i++)
        for (unsigned int j = i + 1; j < tab_size; j++) {
            int distance = POPCOUNT32(tab[i] ^ tab[j]);
            if (distance < best) {
                best = distance;
                if (best == lowest)  // nothing smaller is possible
                    return best;
            }
        }
    return best;
}
\end{lstlisting}

% RTL simulation waveforms (Icarus Verilog), made by sim/run_icarus.bat
\begin{figure*}[!t]
\centering
\figorbox{figures/wave_v1}{\textwidth}{4.4cm}{Waveform of the version~1 core:
  run \code{sim/run\_icarus.bat} (creates \code{figures/wave\_v1.png})}\\[2pt]
\figorbox{figures/wave_v2}{\textwidth}{4.4cm}{Waveform of the version~2 core:
  run \code{sim/run\_icarus.bat} (creates \code{figures/wave\_v2.png})}
\caption{RTL simulation of both HLS cores on main.c's 5,000 words at 100\,MHz
(Icarus Verilog, drawn from the VCD files). \code{phase}: the running step;
\code{rows}: rows (version~1) or blocks (version~2) of pairs started;
\code{RVALID}: words arriving from DDR; \code{latency}: cycles from
\code{ap\_start} to \code{ap\_done}. Next to the names are the final values; the
cursor marks \code{ap\_done}.}
\label{fig:waves}
\end{figure*}

\subsection{The FPGA Cores}
For the FPGA, I rewrote both algorithms in C++ for Vitis HLS with the same
interface. Each core takes up to 8,192 words (what fits on chip) and returns
$-2$ for more. The version~1 core copies the words from DDR into 16 block-RAM
banks, then checks 16 pairs per clock, each with an XOR and an adder tree that
counts the bits, and never stops early. The version~2 core keeps the three steps
with parts that suit hardware: a 4-pass radix sort~\cite{knuth3} instead of
\code{qsort}, a 13-level search tree, built twice, that answers 4 look-ups per
clock, and a pair scan of 64 pairs per clock with an early exit. To feed them,
it keeps six copies of the input on chip.

% ============================================================================
\section{Testing and Benchmarks}

\subsection{Benchmark and Tests}
The same \code{main.c} drives both versions. It runs 6 tests with answers worked
out by hand (a general case, a duplicate, one bit apart, two bits apart, the
largest distance and a single word: 4, 0, 1, 2, 32 and $-1$). Then it times one
call on 5,000 words from the MT19937 generator~\cite{mt19937} with seed 42. Each
build ran 15 times, and I report the median.

The smallest distance in this array is 2. Three pairs reach it, and the first is
words 263 and 3,180 (Table~\ref{tab:pair}). Version~1 checks all 12,497,500
pairs. Version~2 finds no duplicate in step~1 and no one-bit neighbour in
step~2 (160,000 binary searches), then stops in step~3 after 1,283,201 pairs,
about 10\,\% of them.

\begin{table}[!t]
\caption{Closest Pair of the Benchmark Array: Words 263 and 3,180 Differ Only
in Bits 9 and 4}
\label{tab:pair}
\centering\footnotesize
\setlength{\tabcolsep}{3.5pt}
\begin{tabular}{@{}lll@{}}
\toprule
 & \textbf{Hex} & \textbf{Binary, bit 31 \dots{} bit 0} \\
\midrule
word 263   & \code{B1437CE0} & \code{1011 0001 0100 0011 0111 11\textbf{0}0 111\textbf{0} 0000} \\
word 3,180 & \code{B1437EF0} & \code{1011 0001 0100 0011 0111 11\textbf{1}0 111\textbf{1} 0000} \\
XOR        & \code{00000210} & \code{0000 0000 0000 0000 0000 00\textbf{1}0 000\textbf{1} 0000} \\
\bottomrule
\end{tabular}
\end{table}

A sweep program (\code{bench.c}) times both versions on 5 sizes (1,000 to 20,000
words), 7 kinds of input and both builds, 70 measurements in all, on one pinned
core (median of 3 to 9 calls). Besides main.c's array, the inputs use code words
of the extended Hamming code [32,\,26,\,4]~\cite{hamming}, where any two words
differ in at least 4 bits, with one planted pair at distance 0, 1 or 2, or with
none; then the answer is 4 and version~2 cannot stop early.

Both versions return the same answer on every test and on 20,005 random arrays
of 0 to 2,500 words. I measured the heap with Valgrind Massif~\cite{valgrind},
the stack with a painted-stack program (\code{stack2.c}) that runs each call on
a fresh stack filled with a pattern and counts the bytes that changed, and the
instructions inside the call with Callgrind.

\subsection{RTL Simulation}
I also simulated the Verilog that Vitis HLS generated for each core with Icarus
Verilog~\cite{icarus}, using a testbench (\code{tb\_min\_hamming.v}) that acts
like the ARM program: it writes the array address, the size and
\code{ap\_start} over AXI4-Lite~\cite{axi}, serves the core's AXI4 read bursts
from a memory model holding the same 5,000 words, and polls \code{ap\_done}.
Fig.~\ref{fig:waves} shows both runs. The version~1 core copies the words in
(50\,$\mu$s), works through 4,999 rows of pairs and finishes after 895,929
cycles (8.96\,ms at 100\,MHz). The version~2 core loads
(50\,$\mu$s), sorts (421\,$\mu$s), builds its tree (50\,$\mu$s) and runs step~2
(400\,$\mu$s); step~3 then finds a distance-2 pair after 264 blocks, and the
core finishes after 114,981 cycles (1.150\,ms), as on the board. Both return 2.
On the board, version~1 took 9.008\,ms because its timer also covers the cache
flush and the register accesses.

% Vivado block design of version 2 (figures/bd_v2.pdf; a .png also works).
% Version 1: figures/bd_v1.pdf (same blocks, Zybo Z7-20).
\begin{figure*}[!t]
\centering
\figorbox{figures/bd_v2}{0.78\textwidth}{4.6cm}{Block design:
  put the screenshot in \code{figures/bd\_v2.png}}
\caption{Vivado block design of the version~2 system on the ZC702. The version~1
system on the Zybo Z7-20 has the same blocks and connections.}
\label{fig:bd}
\end{figure*}

% ============================================================================
\section{Platforms and FPGA Deployment}

\emph{PC.} An Intel Core i5-1135G7 (Tiger Lake, up to 4.2\,GHz) running Ubuntu
24.04 with GCC 13.3, glibc 2.39 and Valgrind 3.22; calls are timed with
\code{clock\_gettime}. I used two builds: scalar (\code{gcc -O2 -mpopcnt
-std=c11}) and AVX2 (\code{gcc -O3 -march=native -std=c11}). In the AVX2 build,
GCC vectorizes the inner loop of version~1, while the early-exit loops of
version~2 stay scalar.

\emph{Boards.} Version~1 ran on a Zybo Z7-20 and version~2 on a ZC702, both
with the XC7Z020-1 chip and the same clocks: 666.67\,MHz for the ARM, 533\,MHz
for the DDR3 and 100\,MHz for the FPGA logic. The ARM code is bare metal
(standalone BSP, \code{arm-none-eabi-gcc -O2}, newlib), timed with the Cortex-A9
global timer; the tools were Vitis HLS, Vivado and Vitis 2021.1.

\emph{Block design.} Each core is an IP block next to the Zynq processing
system~\cite{ug585} (Fig.~\ref{fig:bd}). The ARM drives the core's AXI4-Lite
registers at \code{0x43C0\_0000} through \code{M\_AXI\_GP0} and an AXI
Interconnect: start and done at 0x00, the result at 0x10, the array address at
0x18 and the size at 0x20. The core reads the words itself: its AXI4 master,
\code{m\_axi\_gmem}, goes through a second interconnect into \code{S\_AXI\_HP0}
and on to DDR, in bursts. The program flushes the array from the data cache,
writes the registers, starts the core and polls \code{ap\_done}. Since
version~2 calls \code{malloc}, its board build reserves a 1\,MB heap
(\code{-Wl,-{}-defsym=\_HEAP\_SIZE=0x100000}). With the BSP default of 8\,KB,
\code{malloc} would fail near $N=2{,}048$ and version~2 would lose steps 1 and 2
(but still answer correctly).

\emph{Board tests.} Version~1 passed 7 tests: main.c's six on both the core and
the ARM, and an array of 8,193 words that the core rejects with $-2$.
Version~2 passed 8 tests (the same six, an empty array and 8,193 words) and 300
random arrays checked against the ARM version with no mismatch
(Listing~\ref{lst:uart}). Both ARM times in Section~\ref{sec:results} come from
the version~2 program, which contains both versions; version~1's own program
measured 2,737.1\,ms.

\begin{lstlisting}[float=t,language={},numbers=none,xleftmargin=0pt,basicstyle=\ttfamily\fontsize{6}{7}\selectfont,caption={UART output of the version 2 board program (excerpt)},label={lst:uart}]
-- Benchmark: N = 5000 random words (MT19937, seed 42) --
  ARM version 1 : distance 2, 2699.550 ms
  ARM version 2 : distance 2, 72.437 ms  (fastest of 5)
  PL core       : distance 2, 1.150 ms .. 1.151 ms  (10 runs, decided in step 3)
-- Random regression: 300 arrays, PL core vs ARM version 2 --
  0 mismatches; decided by step 1: 52, step 2: 69, step 3: 179
ALL TESTS PASSED
\end{lstlisting}

% ============================================================================
\section{Results and Comparison}
\label{sec:results}
Table~\ref{tab:score} lists all the results. The estimates come from Vivado's
power report and the HLS cycle model and are marked as such; the rest were
measured.

\begin{table}[!t]
\caption{Scorecard (One Call, 5,000 Words, Unless Noted)}
\label{tab:score}
\centering\footnotesize
\setlength{\tabcolsep}{2.5pt}
\begin{tabularx}{\columnwidth}{@{}>{\raggedright\arraybackslash}Xrrcr@{}}
\toprule
\textbf{Measure} & \textbf{Version 1} & \textbf{Version 2} & \textbf{Better} & \textbf{By} \\
\midrule
\multicolumn{5}{@{}l}{\emph{Processing time}} \\
PC, scalar build             & 303.3\,ms   & 5.90\,ms  & V2 & 51$\times$ \\
PC, AVX2 build               & 34.9\,ms    & 5.98\,ms  & V2 & 5.8$\times$ \\
ARM Cortex-A9                & 2,699.6\,ms & 72.4\,ms  & V2 & 37$\times$ \\
FPGA core, 100\,MHz          & 9.008\,ms   & 1.150\,ms & V2 & 7.8$\times$ \\
PC sweep, 70 cases           & --          & 70 wins   & V2 & 1.34--3,222$\times$ \\
No pair under 4 bits, PC     & 303.5\,ms   & 10.8\,ms  & V2 & 28$\times$ \\
PC, duplicate, $N=20{,}000$  & 5,030\,ms   & 1.67\,ms  & V2 & 3,007$\times$ \\
Instructions, PC scalar      & 3,336.9\,M  & 33.0\,M   & V2 & 101$\times$ \\
\addlinespace[3pt]
\multicolumn{5}{@{}l}{\emph{Memory size}} \\
Extra heap, PC               & 0           & $8N$\,B   & V1 & -- \\
Stack during the call, PC    & ${<}\,128$\,B & 2,712\,B & V1 & ${>}\,21\times$ \\
Code, PC                     & 222\,B      & 606\,B + libc & V1 & 2.7$\times$ \\
Extra heap, ARM              & 0           & $4N$\,B   & V1 & -- \\
Code and data, ARM           & 128\,B      & 6,810\,B  & V1 & 53$\times$ \\
FPGA LUTs                    & 2,557       & 11,898    & V1 & 4.7$\times$ \\
FPGA flip-flops              & 3,328       & 8,314     & V1 & 2.5$\times$ \\
FPGA block RAM tiles         & 12          & 52.5      & V1 & 4.4$\times$ \\
On-chip copies of the input  & 32\,KB      & 192\,KB   & V1 & 6.0$\times$ \\
\addlinespace[3pt]
\multicolumn{5}{@{}l}{\emph{Time and size together (FPGA cores, estimates)}} \\
Dynamic power                & 0.090\,W    & 0.331\,W  & V1 & 3.7$\times$ \\
Energy per call              & 0.81\,mJ    & 0.38\,mJ  & V2 & 2.1$\times$ \\
Same, no pair under 4 bits   & 0.81\,mJ    & 1.13\,mJ  & V1 & 1.4$\times$ \\
\addlinespace[3pt]
Correctness, all tests       & pass        & pass      & tie & -- \\
\bottomrule
\end{tabularx}
\end{table}

\emph{Processing time: version~2 wins.}
\begin{itemize}
\item It is faster everywhere: 51$\times$ on the PC, 5.8$\times$ with AVX2,
  37$\times$ on the ARM and 7.8$\times$ on the FPGA, with 101$\times$ fewer
  instructions.
\item The gap follows the cost of one pair. AVX2 brings version~1 down to 37
  instructions, and on the ARM, which has no popcount instruction, version~2
  calls libgcc's \code{\_\_popcountsi2} and needs about 21 against about 201.
\item It wins all 70 sweep cases (1.34$\times$ to 3,222$\times$), and even with
  no pair closer than 4 bits, where it cannot stop early, it stays 28$\times$
  (scalar) and 2.9$\times$ (AVX2) faster.
\end{itemize}

\emph{Memory size: version~1 wins.}
\begin{itemize}
\item PC: no extra heap against $8N$ bytes, under 128\,B of stack against
  2,712\,B, and 222\,B of code against 606\,B plus \code{qsort}, \code{malloc}
  and \code{memcpy}.
\item ARM: no extra heap against $4N$ bytes, and 53$\times$ less code and data
  with the library code counted (128 against 6,810\,B).
\item FPGA: 4.7$\times$ fewer LUTs, 2.5$\times$ fewer flip-flops and 4.4$\times$
  less block RAM. The whole design uses 6.2\,\% of the LUTs instead of
  23.8\,\%, with more slack at 100\,MHz (+2.952 against +0.672\,ns).
\end{itemize}

\emph{Time and size together.} Version~2's core draws more power (0.331 against
0.090\,W, estimate) but finishes 7.8 times sooner, so it needs 2.1 times less
energy per call (0.38 against 0.81\,mJ) and less LUT-time (13,683 against
23,033\,LUT$\cdot$ms). With no pair closer than 4 bits this reverses (1.13\,mJ
and 40,493\,LUT$\cdot$ms, estimate), and version~1 is the more efficient
circuit, though 2.6 times slower.

\emph{Correctness is a tie:} the same answers on every test, on the boards and
on 20,005 random arrays.

% ============================================================================
\section{Conclusion}
The two versions sit at the two ends of the time--memory trade-off. Version~2
wins every time measure on every platform, and its lead grows with $N$, since
its sort and neighbour search grow as $N\log N$ while the pair scan that it
usually cuts short grows as $N^2$; its memory cost stays a fixed multiple of the
input. Version~1 wins every memory measure: no heap, a tiny stack, 128\,B of
code on the ARM and the smallest FPGA core.

I would use version~2 when time matters and the target can spare $O(N)$ memory:
a PC or application processor, the Zynq ARM with a raised heap, or an FPGA with
block RAM to spare. Version~1 fits targets with no heap or a few kilobytes of
RAM, where code size counts, and FPGA designs where area or energy matter more
than latency and close pairs are rare. Version~2's memory could also come down
(not built or measured): sorting \code{tab[]} in place, if the caller allows it,
removes the extra heap, and an in-place sort such as heapsort removes glibc's
$4N$-byte buffer and most of the stack.

My answer to Problem~1 is version~2. Version~1 remains the right choice when
memory, not time, is the limit.

% ============================================================================
\hbadness=10000
\begin{thebibliography}{9}
\bibitem{hamming} R.~W. Hamming, ``Error detecting and error correcting codes,''
  \emph{Bell Syst. Tech. J.}, vol.~29, no.~2, pp.~147--160, Apr. 1950.
\bibitem{knuth3} D.~E. Knuth, \emph{The Art of Computer Programming, Vol.~3:
  Sorting and Searching}, 2nd~ed. Reading, MA, USA: Addison-Wesley, 1998.
\bibitem{mt19937} M.~Matsumoto and T.~Nishimura, ``Mersenne twister: A
  623-dimensionally equidistributed uniform pseudo-random number generator,''
  \emph{ACM Trans. Model. Comput. Simul.}, vol.~8, no.~1, pp.~3--30, Jan. 1998.
\bibitem{valgrind} N.~Nethercote and J.~Seward, ``Valgrind: A framework for
  heavyweight dynamic binary instrumentation,'' in \emph{Proc. ACM SIGPLAN Conf.
  Program. Lang. Design Implement. (PLDI)}, 2007, pp.~89--100.
\bibitem{ug1399} Xilinx, \emph{Vitis High-Level Synthesis User Guide}, UG1399,
  2021.
\bibitem{ug585} Xilinx, \emph{Zynq-7000 SoC Technical Reference Manual}, UG585.
\bibitem{axi} Arm, \emph{AMBA AXI and ACE Protocol Specification}, IHI 0022.
\bibitem{icarus} S.~Williams, \emph{Icarus Verilog}. [Online]. Available:
  \url{https://github.com/steveicarus/iverilog}
\end{thebibliography}

\end{document}
